\subsection{Study QJ: Text-to-Image Scenes Scored by the Frozen Gate}\label{sec:moved-qj}
Study QI edited simulator frames and could not say how the gate behaves on scenes that never came from the simulator. Study QJ (registered, OpenTimestamps-stamped, Amendment 1 before any study image) asks that with Qwen-Image-2.1-Turbo used as a plain text-to-image model: no simulator frame, no mask, no input image; six scene templates (two supplies, three lightings), twenty seeds each, \QJN{} images per mode: clean (N0), a small partly hidden fire behind the crates (NF, the paper's hard class) and a large fire (NB, a positive control). The frozen gate scores each image (one view, $\tau=0.088$ from XD, nothing refitted).
\emph{Claims.} QJ1: at most 6 of 120 clean scenes reach $\tau$: \QJFalse{} of \QJN{} (upper \QJFalseUpper), \QJClaimQJOne. QJ2: at most 10\% of the large fires are missed: \QJLargeMiss{} of \QJLargeN{} (\QJLargeRate, upper \QJLargeUpper), \QJClaimQJTwo. QJ3: at most 20\% of the small fires are missed: \QJSmallMiss{} of \QJSmallN{} (\QJSmallRate, upper \QJSmallUpper), \QJClaimQJThree. A by-eye sample of 40 images per fire mode, classified before any score, shows a fire in \QJLargeSeen{} of \QJLargeSample{} (NB) and \QJSmallSeen{} of \QJSmallSample{} (NF) images; the registered stop rule needs at least 34.
\emph{Reading.} All three claims hold, and none is close to its line: no clean scene is flagged, and no large or small generated fire is missed, in any of the six templates (upper bound \QJLargeUpper{} over the 120 images of a mode). The stop rule did not trigger: the generator drew a fire in nearly every sampled image, which is the opposite of QI-G, where 68\% of the masked edits carried no flame colour. Read together with QI, the picture is narrower than either alone: the gate recognises photographic flames and smoke when the image contains them, and it keeps a photographic clean scene clean; what it loses in QI1 is the simulator's grey smoke discs once they sit on a photographic background. The hard class of the simulation is therefore hard partly because of how its smoke is rendered, not only because of where it stands. \textbf{This is not a result about the hard class.} The prompt asks for a small fire partly hidden behind the crates, but in the sampled images the flames are mostly plainly visible above or beside the crates, so the generated ``small hidden fire'' is easier than the simulator's occluded one; QJ3 holds for what was generated, and the occlusion geometry of the hard class is not reproduced. Limits: generated images are not photographs and no real site was seen; the generator and the gate are both Qwen-family models; one view, with the three-view threshold of XD; one checkpoint and one wording per mode; the by-eye reader is a single reader; the claim is about the gate's recognition of generated fire, not about missions.
